\chapter{Number Theory - Fundamentals and Advanced Topics}
\label{ch:number-theory}

\section{Divisibility and GCD}
\label{sec:divisibility-gcd}

Integers $\Z$ form an integral domain with the usual operations $+$, $-$, $\cdot$. For $a, b \in \Z$ with $b > 0$, we say $a$ is \textbf{divisible} by $b$ if there exists $k \in \Z$ such that $a = bk$. We write $b \mid a$.

\begin{definition}[Greatest Common Divisor]
\label{def:gcd}
For $a, b \in \Z \setminus \{0\}$, the \textbf{greatest common divisor} $\gcd(a, b)$ is the largest positive integer dividing both $a$ and $b$.
\end{definition}

\begin{theorem}[Euclid's Lemma]
\label{thm:euclid-lemma}
Let $a, b, c \in \Z$ with $a \neq 0$. If $a \mid bc$ and $\gcd(a, b) = 1$, then $a \mid c$.
\end{theorem}

\begin{proof}
By Bézout's identity, since $\gcd(a, b) = 1$, there exist $x, y \in \Z$ such that $ax + by = 1$. Multiply by $c$: $acx + bc y = c$. Since $a \mid bc$, let $bc = ka$. Then $acx + ka y = c$, so $a(cx + ky) = c$. Thus $a \mid c$.
\end{proof}

\begin{definition}[Euclid's Algorithm]
\label{def:euclid-algorithm}
To compute $\gcd(a, b)$ with $a > b > 0$:
\begin{enumerate}
\item $r_0 = a, r_1 = b$
\item For $n \geq 1$, $r_{n+2} = r_n \bmod r_{n+1}$ (the remainder of $r_n$ divided by $r_{n+1}$)
\item Stop when $r_{k+1} = 0$. Then $\gcd(a, b) = r_k$.
\end{enumerate}
\end{definition}

\begin{theorem}[Euclid's Lemma for GCD]
\label{thm:euclid-lemma-gcd}
$\gcd(a, b) = r_{k-1}$ where $r_k = 0$ in Euclid's algorithm.
\end{theorem}

\begin{proof}
By construction, $r_{k-2} = q_k r_{k-1} + r_k$ and $r_k = 0$, so $r_{k-2} = q_k r_{k-1}$. Thus $r_{k-1} \mid r_{k-2}$. Also $r_{k-1} \mid r_{k-3}$ by the same argument. By induction, $r_{k-1} \mid r_i$ for all $i \leq k-1$. Thus $r_{k-1} \mid a$ and $r_{k-1} \mid b$. Since $r_{k-2} = q_k r_{k-1}$, any common divisor of $a$ and $b$ divides $r_{k-2}$, hence divides $r_{k-1}$. Thus $r_{k-1} = \gcd(a, b)$.
\end{proof}

\section{Prime Numbers}
\label{sec:primes}

A positive integer $p > 1$ is \textbf{prime} if its only divisors are $1$ and $p$.

\begin{theorem}[Fundamental Theorem of Arithmetic]
\label{thm:fta}
Every integer $n > 1$ can be written uniquely (up to order) as a product of primes: $n = p_1 p_2 \cdots p_k$ where $p_1 \leq p_2 \leq \cdots \leq p_k$.
\end{theorem}

\begin{proof}
Existence: Proceed by strong induction. For $n = 2$, prime. Assume true for all $m < n$. If $n$ is prime, done. If composite, $n = ab$ with $1 < a, b < n$. By induction, $a = p_1 \cdots p_i$ and $b = q_1 \cdots q_j$. Thus $n = p_1 \cdots p_i q_1 \cdots q_j$ is a prime factorization.
Uniqueness: Suppose $n = p_1 \cdots p_k = q_1 \cdots q_m$. Then $p_1 \mid q_1 \cdots q_m$. Since $p_1$ is prime, $p_1 \mid q_j$ for some $j$. Relabel so $p_1 = q_1$. Divide and proceed by induction.
\end{proof}

\section{Coprime and Relatively Prime}
\label{sec:coprime}

Integers $a, b$ are \textbf{coprime} (or \textbf{relatively prime}) if $\gcd(a, b) = 1$.

\begin{theorem}[Bézout's Identity]
\label{thm:bézout}
For any $a, b \in \Z$, there exist $x, y \in \Z$ such that $ax + by = \gcd(a, b)$.
\end{theorem}

\begin{proof}
Let $d = \gcd(a, b)$. Consider the set $S = \{ax + by : x, y \in \Z, ax + by > 0\}$. By the well-ordering principle, $S$ has a minimum element $d$. Let $ax_0 + by_0 = d$. If $d \neq \gcd(a, b)$, let $d' = \gcd(a, b)$. Then $d' \mid a$ and $d' \mid b$, so $d' \mid (ax_0 + by_0) = d$. Thus $d \leq d'$. But also $\gcd(a, b) \mid \gcd(a, b) = d$, so $d' \leq d$. Thus $d = d'$.
\end{proof}

\section{Euclid's Proof of Infinite Primes}
\label{sec:euclid-primes}

\begin{theorem}[Euclid's Theorem]
\label{thm:euclid-primes}
There are infinitely many primes.
\end{theorem}

\begin{proof}
Assume there are finitely many primes: $p_1, p_2, \ldots, p_k$. Consider $N = p_1 p_2 \cdots p_k + 1$. Then $N > 1$, so $N$ has a prime factor $q$. If $q = p_i$ for some $i$, then $q \mid N$ and $q \mid p_1 \cdots p_k$, so $q \mid 1$, impossible. Thus $q$ is a new prime, contradiction.
\end{proof}

\section{Congruences and Modular Arithmetic}
\label{sec:congruences}

For $a, b \in \Z$ and $n \in \Z_{\geq 1}$, we write $a \equiv b \pmod n$ if $n \mid (a - b)$.

\begin{theorem}[Congruence Properties]
\label{thm:congruence-properties}
The following hold for $a, b, c, d \in \Z$ and $n \in \Z_{\geq 1}$:
\begin{enumerate}
\item If $a \equiv b \pmod n$, then $a = b + kn$ for some $k \in \Z$.
\item $a \equiv b \pmod n \iff a - b = kn$ for some $k \in \Z$.
\item $a \equiv b \pmod n \text{ and } a \equiv c \pmod n \implies b \equiv c \pmod n$.
\item $a \equiv b \pmod n \text{ and } c \equiv d \pmod n \implies ac \equiv bd \pmod n$.
\item $a \equiv b \pmod n \text{ and } c \equiv d \pmod n \implies a + c \equiv b + d \pmod n$.
\item $a \equiv b \pmod n \implies a^k \equiv b^k \pmod n$ for $k \in \Z_{\geq 1}$.
\end{enumerate}
\end{theorem}

\begin{proof}
(1) and (2) are equivalent definitions. (3) follows from (2): if $a = b + kn_1$ and $a = c + kn_2$, then $b + kn_1 = c + kn_2$, so $b - c = k(n_2 - n_1)$. (4) and (5) follow from $a = b + kn_1$ and $c = d + kn_2$: $ac = bd + bkn_2 + dkn_1 + k^2n_1n_2 = bd + kn_2(b + dkn_1/k + kn_1n_2/n_2)$. (6) follows by induction.
\end{proof}

\section{Chinese Remainder Theorem}
\label{sec:chinese-remainder}

\begin{theorem}[Chinese Remainder Theorem]
\label{thm:chinese-remainder}
Let $n_1, n_2, \ldots, n_k$ be pairwise coprime positive integers. For any $a_1, a_2, \ldots, a_k \in \Z$, there exists a unique $x \in \Z$ modulo $n_1 n_2 \cdots n_k$ such that $x \equiv a_i \pmod{n_i}$ for all $i$.
\end{theorem}

\begin{proof}
Existence: Let $N = n_1 \cdots n_k$ and $N_i = N/n_i$. Since $\gcd(N_i, n_i) = 1$, by Bézout's identity, there exist $u_i, v_i \in \Z$ such that $N_i u_i + n_i v_i = 1$. Let $x = \sum_{i=1}^k a_i N_i u_i$. Then $x \equiv a_i N_i u_i \equiv a_i \cdot 1 \equiv a_i \pmod{n_i}$. Uniqueness: If $x \equiv y \pmod{N}$, then $x - y = mN = m n_1 \cdots n_k$, so $n_i \mid (x - y)$ for all $i$.
\end{proof}

\section{Fermat's Little Theorem}
\label{sec:fermat-little}

\begin{theorem}[Fermat's Little Theorem]
\label{thm:fermat-little}
Let $p$ be prime and $a \in \Z$. If $p \nmid a$, then $a^{p-1} \equiv 1 \pmod p$.
\end{theorem}

\begin{proof}
Consider the set $S = \{1, 2, \ldots, p-1\}$. Multiplication by $a$ permutes $S$ modulo $p$ (since $p \nmid a$). Thus $\prod_{x \in S} x \equiv \prod_{x \in S} ax \pmod p$. The LHS is $(p-1)!$ and the RHS is $a^{p-1}(p-1)!$. Thus $a^{p-1} \equiv 1 \pmod p$.
\end{proof}

\section{Euclid's Lemma for Primes}
\label{sec:euclid-primes-lemma}

\begin{theorem}[Euclid's Lemma for Primes]
\label{thm:euclid-primes-lemma}
If $p$ is prime and $p \mid ab$, then $p \mid a$ or $p \mid b$.
\end{theorem}

\begin{proof}
If $p \mid a$, done. If $\gcd(p, a) = 1$, then by Euclid's lemma (for groups), $p \mid b$.
\end{proof}

\section{Wilson's Theorem}
\label{sec:wilson}

\begin{theorem}[Wilson's Theorem]
\label{thm:wilson}
A positive integer $p$ is prime if and only if $(p-1)! \equiv -1 \pmod p$.
\end{theorem}

\begin{proof}
(Only if) Let $p$ be prime. For $a = 1, 2, \ldots, p-1$, $a \not\equiv 1/a \pmod p$ (otherwise $a^2 \equiv 1 \pmod p$, so $a \equiv 1$ or $a \equiv p-1$). Pairing elements $a$ with $a^{-1} \not\equiv a$, we get $(p-3)/2$ pairs of $(a, a^{-1})$ with $a \not\equiv a^{-1}$, and two fixed points $1$ and $-1$. Thus $(p-1)! \equiv 1 \cdot (-1) \cdot \prod a \cdot a^{-1} \equiv -1 \pmod p$.
(If) If $(p-1)! \equiv -1 \pmod p$, then $p \mid (p-1)! + 1$. If $p$ is composite, let $p = ab$ with $1 < a, b < p$. Then $a, b \leq p-1$, so $ab \leq (p-1)^2 < p(p-1)$. But $(p-1)!$ contains $a$ and $b$ as factors, so $ab \mid (p-1)!$. Thus $p \mid (p-1)!$, contradiction.
\end{proof}

\section{Exercises}
\label{sec:number-theory-exercises}

\begin{exercise}[GCD of Coprime Numbers]
\label{ex:gcd-coprime}
Show that $\gcd(a, b) = 1 \iff \gcd(a, bc) = 1 \text{ or } \gcd(b, ac) = 1$.
\begin{proof}
If $\gcd(a, b) = 1$, then $\gcd(a, bc) = 1$ or $\gcd(b, ac) = 1$. Conversely, if $\gcd(a, b) > 1$, let $p = \gcd(a, b)$. Then $p \mid a$ and $p \mid b$, so $p \mid ac$ and $p \mid bc$, so $\gcd(a, bc) \geq p > 1$.
\end{proof}
\end{exercise}

\begin{exercise}[Congruence and Primality Test]
\label{ex:congruence-primality}
Use Fermat's little theorem to determine if $n = 211$ is prime.
\begin{proof}
Compute $210! \pmod{211}$. If $211$ is prime, then $210! \equiv -1 \pmod{211}$ by Wilson's theorem. If $211$ is composite, let $p \mid 211$. Then $p \leq \sqrt{211} < 15$. If $p \mid 211!$, then $p \mid 210!$. But $210! \equiv -1 \pmod{211}$, so $p \nmid 210!$. Contradiction. Thus $211$ is prime.
\end{proof}
\end{exercise}
